% GENERATED FILE. Edit canonical lessons, then run npm run generate:books.
% canonical-source-hash: fnv1a64:4a7308622da4446e
% canonical-lessons: RU-C455-rasplachivatsya, RU-C455-oplata, RU-C455-kartochka, RU-C455-razmenivat, RU-C455-obschityvat

\chapter{Paying at the till}
\label{ch:ru-paying-at-the-till}
\hlchaptermodality{455}

\begin{chapteropening}
\textbf{By the end of this chapter:} \emph{I can pay: settling up, the payment, a bank card, changing money and being shortchanged.}

Complete the last lesson of chapter 455: I can pay: settling up, the payment, a bank card, changing money and being shortchanged.
\end{chapteropening}

\section[\texorpdfstring{\ru{расплачиваться} (raspláchivatsya)}{расплачиваться (raspláchivatsya)}]{\ru{расплачиваться} (raspláchivatsya) — to settle up, to pay (the bill)}

\label{lesson:RU-C455-rasplachivatsya}



\begin{quote}
Before the new one: say the Russian for to haggle, then the Russian for to overpay.
\end{quote}

\subsection*{You'll want to know: \ru{расплачиваться}}
\textbf{\ru{расплачиваться}} (\emph{raspláchivatsya}) — \textquotedblleft{}to settle up, to pay (the bill)\textquotedblright{}.

\begin{quote}
\ru{Мы} \ru{расплачиваемся} \ru{и} \ru{уходим}.
\end{quote}

\begin{quote}
\emph{My raspláchivayemsya i ukhódim.}
\end{quote}

\begin{quote}
We settle up and leave.
\end{quote}

Imperfective; the perfective is \emph{\ru{расплатиться}}. Takes \emph{\ru{за}} + accusative: \emph{\ru{расплатиться} \ru{за} \ru{ужин}}.

\begin{cousinweb}
\emph{\ru{расплачиваться}} is built from \emph{\ru{рас}-}, \emph{-\ru{плачивать}}, from \emph{\ru{платить}} 'to pay', and \emph{-\ru{ся}}. EtymDB takes \emph{\ru{платить}} from Proto-Slavic \emph{platiti}.
\end{cousinweb}

\begin{grammarlens}[title={What you've built}]
One more word for telling, explaining and describing at B1.
\end{grammarlens}

\subsection*{Your turn}
\begin{itemize}
\raggedright
  \item \emph{Say it:} \emph{raspláchivatsya}
  \item \emph{Say it:} \emph{raspláchivatsya}, once more
  \item \emph{Say it:} a sentence of your own with \emph{raspláchivatsya}
  \item \emph{Recall:} say the Russian for to haggle, then the Russian for to overpay, then say \emph{raspláchivatsya} again
\end{itemize}

\begin{tcolorbox}[breakable,colback=teal!4,colframe=teal!35!black,title={Before you move on}]
What does \ru{расплачиваться} mean? (\textquotedblleft{}To settle up, to pay (the bill)\textquotedblright{}.) Say it once more.
\end{tcolorbox}
\section[\texorpdfstring{\ru{оплата} (opláta)}{оплата (opláta)}]{\ru{оплата} (opláta) — payment}

\label{lesson:RU-C455-oplata}



\begin{quote}
Before the new one: say the Russian for to overpay, then the Russian for to settle up.
\end{quote}

\subsection*{You'll want to know: \ru{оплата}}
\textbf{\ru{оплата}} (\emph{opláta}) — \textquotedblleft{}payment\textquotedblright{}.

\begin{quote}
\ru{Оплата} \ru{только} \ru{наличными}.
\end{quote}

\begin{quote}
\emph{Opláta tólko nalíchnymi.}
\end{quote}

\begin{quote}
Payment is in cash only.
\end{quote}

Feminine. \emph{\ru{Оплата} \ru{картой}} — payment by card.

\begin{cousinweb}
\emph{\ru{оплата}} is built from \emph{\ru{о}-} and \emph{\ru{плата}} 'pay', the root of \emph{\ru{платить}} 'to pay', which EtymDB takes from Proto-Slavic \emph{platiti}.
\end{cousinweb}

\begin{grammarlens}[title={What you've built}]
One more word for telling, explaining and describing at B1.
\end{grammarlens}

\subsection*{Your turn}
\begin{itemize}
\raggedright
  \item \emph{Say it:} \emph{opláta}
  \item \emph{Say it:} \emph{opláta}, once more
  \item \emph{Say it:} a sentence of your own with \emph{opláta}
  \item \emph{Recall:} say the Russian for to overpay, then the Russian for to settle up, then say \emph{opláta} again
\end{itemize}

\begin{tcolorbox}[breakable,colback=teal!4,colframe=teal!35!black,title={Before you move on}]
What does \ru{оплата} mean? (\textquotedblleft{}Payment\textquotedblright{}.) Say it once more.
\end{tcolorbox}
\section[\texorpdfstring{\ru{карточка} (kártochka)}{карточка (kártochka)}]{\ru{карточка} (kártochka) — a bank card; a small card}

\label{lesson:RU-C455-kartochka}



\begin{quote}
Before the new one: say the Russian for to settle up, then the Russian for payment.
\end{quote}

\subsection*{You'll want to know: \ru{карточка}}
\textbf{\ru{карточка}} (\emph{kártochka}) — \textquotedblleft{}a bank card; a small card\textquotedblright{}.

\begin{quote}
\ru{Можно} \ru{заплатить} \ru{карточкой}?
\end{quote}

\begin{quote}
\emph{Mózhno zaplatít kártochkoy?}
\end{quote}

\begin{quote}
Can I pay by card?
\end{quote}

Feminine; a diminutive of \emph{\ru{карта}}. In speech \emph{\ru{карточка}} is a bank card.

\begin{cousinweb}
\emph{\ru{карточка}} is a diminutive of \emph{\ru{карта}} 'a card', which EtymDB traces through Polish \emph{karta} and German \emph{Karte}; the source is Latin \emph{charta} 'paper' (Lewis \& Short).
\end{cousinweb}

\begin{grammarlens}[title={What you've built}]
One more word for telling, explaining and describing at B1.
\end{grammarlens}

\subsection*{Your turn}
\begin{itemize}
\raggedright
  \item \emph{Say it:} \emph{kártochka}
  \item \emph{Say it:} \emph{kártochka}, once more
  \item \emph{Say it:} a sentence of your own with \emph{kártochka}
  \item \emph{Recall:} say the Russian for to settle up, then the Russian for payment, then say \emph{kártochka} again
\end{itemize}

\begin{tcolorbox}[breakable,colback=teal!4,colframe=teal!35!black,title={Before you move on}]
What does \ru{карточка} mean? (\textquotedblleft{}A bank card; a small card\textquotedblright{}.) Say it once more.
\end{tcolorbox}
\section[\texorpdfstring{\ru{разменивать} (razménivat)}{разменивать (razménivat)}]{\ru{разменивать} (razménivat) — to break (a note), to change (money into smaller notes)}

\label{lesson:RU-C455-razmenivat}



\begin{quote}
Before the new one: say the Russian for payment, then the Russian for a bank card.
\end{quote}

\subsection*{You'll want to know: \ru{разменивать}}
\textbf{\ru{разменивать}} (\emph{razménivat}) — \textquotedblleft{}to break (a note), to change (money into smaller notes)\textquotedblright{}.

\begin{quote}
\ru{В} \ru{кассе} \ru{не} \ru{разменивают} \ru{крупные} \ru{деньги}.
\end{quote}

\begin{quote}
\emph{V kásse ne razménivayut krúpnyye déngi.}
\end{quote}

\begin{quote}
They do not change large notes at the till.
\end{quote}

Imperfective; perfective \emph{\ru{разменять}}. \emph{\ru{Разменять} \ru{тысячу}} — to break a thousand-rouble note.

\begin{cousinweb}
\emph{\ru{разменивать}} is built from \emph{\ru{раз}-} and \emph{-\ru{менивать}}, from \emph{\ru{менять}} 'to change', on \emph{\ru{мена}} 'an exchange', which is inherited from Proto-Slavic \emph{měna} (EtymDB).
\end{cousinweb}

\begin{grammarlens}[title={What you've built}]
One more word for telling, explaining and describing at B1.
\end{grammarlens}

\subsection*{Your turn}
\begin{itemize}
\raggedright
  \item \emph{Say it:} \emph{razménivat}
  \item \emph{Say it:} \emph{razménivat}, once more
  \item \emph{Say it:} a sentence of your own with \emph{razménivat}
  \item \emph{Recall:} say the Russian for payment, then the Russian for a bank card, then say \emph{razménivat} again
\end{itemize}

\begin{tcolorbox}[breakable,colback=teal!4,colframe=teal!35!black,title={Before you move on}]
What does \ru{разменивать} mean? (\textquotedblleft{}To break (a note), to change (money into smaller notes)\textquotedblright{}.) Say it once more.
\end{tcolorbox}
\section[\texorpdfstring{\ru{обсчитывать} (obschítyvat)}{обсчитывать (obschítyvat)}]{\ru{обсчитывать} (obschítyvat) — to shortchange}

\label{lesson:RU-C455-obschityvat}



\begin{quote}
Before the new one: say the Russian for a bank card, then the Russian for to break.
\end{quote}

\subsection*{You'll want to know: \ru{обсчитывать}}
\textbf{\ru{обсчитывать}} (\emph{obschítyvat}) — \textquotedblleft{}to shortchange\textquotedblright{}.

\begin{quote}
\ru{Этот} \ru{продавец} \ru{часто} \ru{обсчитывает} \ru{покупателей}.
\end{quote}

\begin{quote}
\emph{Étot prodavéts chásto obschítyvayet pokupáteley.}
\end{quote}

\begin{quote}
This seller often shortchanges customers.
\end{quote}

Imperfective; perfective \emph{\ru{обсчитать}}.

\begin{cousinweb}
\emph{\ru{обсчитывать}} is built from \emph{\ru{об}-} and \emph{-\ru{считывать}}, from \emph{\ru{считать}} 'to count': to count to someone's loss. EtymDB lists \emph{\ru{считать}} without a deeper source.
\end{cousinweb}

\begin{grammarlens}[title={What you've built}]
One more word for telling, explaining and describing at B1.
\end{grammarlens}

\subsection*{Your turn}
\begin{itemize}
\raggedright
  \item \emph{Say it:} \emph{obschítyvat}
  \item \emph{Say it:} \emph{obschítyvat}, once more
  \item \emph{Say it:} a sentence of your own with \emph{obschítyvat}
  \item \emph{Recall:} say the Russian for a bank card, then the Russian for to break, then say \emph{obschítyvat} again
\end{itemize}

\begin{tcolorbox}[breakable,colback=teal!4,colframe=teal!35!black,title={Before you move on}]
What does \ru{обсчитывать} mean? (\textquotedblleft{}To shortchange\textquotedblright{}.) Say it once more.
\end{tcolorbox}
